
% 08. Sexual Distance and the Island Species: A Predator Trapped on an Evolutionary Island.tex

\section{섹슈얼 디스턴스와 아일랜드 종: 진화의 섬에 갇힌 맹수}
우리는 흔히 생태계 최정점에 서 있는 사자나 호랑이 같은 맹수들을 동물의 왕이라 부르며 그들의 독보적인 지위를 찬미한다. 그러나 디스턴스(DISTANCE)의 서늘한 렌즈로 생태계를 해부해 보면, 정점에 선 맹수일수록 그들 주변에 생물학적으로 유사한 종(근접종)이 거의 남아 있지 않다는 섬뜩한 역학적 진실과 마주하게 된다. 그들의 독보적 지위는 평화로운 진화의 결과가 아니라, 자신과 에너지를 다투는 유사종들을 모조리 도태시키고 멸종시킨 맹렬한 경쟁의 결과물이기 때문이다.

이 잔혹하고도 필연적인 우주적 평활화의 궤적 속에서, 인류의 진화사를 돌아보자. 우리는 흔히 호모 사피엔스(인류)가 탁월한 지능을 가졌기에 진화의 정점에 우뚝 섰다고 믿는다. 그러나 인류의 진정한 특수성은 숭고한 지능에 있는 것이 아니다. 인류는 지구라는 생태계의 가장 무자비한 맹수로서, 네안데르탈인, 데니소바인 등 자신과 유전자를 교류할 수 있었던 수많은 유사 인류들을 모조리 멸종시키며 진화의 트랙 위를 홀로 걸어왔다.
그 결과 인류는 생물학적 역사상 전대미문의 기괴한 상태에 도달하게 된다. 주변의 모든 유사종을 잃어버림으로써, 이 지구상에 존재하는 그 어떤 생명체와도 유전자를 섞을 수 없는 완벽한 고립 상태에 빠진 것이다. 물고기나 식물들은 주변에 수많은 근접종들이 살아 숨 쉬고 있어 끊임없이 교배하며 방대한 유전자 조합을 만들어내지만, 인류는 유전자를 교환할 외부의 통로가 완벽하게 막혀버렸다.

디스턴스 에세이에서는 이를 일컬어 가장 극단적인 '섹슈얼 디스턴스(Sexual Distance)'의 형성이라 명명한다. 생명은 끊임없이 에너지를 섞기 위해 다가가고 교배(연결)하지만, 역설적으로 그 맹렬한 연결의 결과로 주변 종들을 멸종시키며 영원히 교배할 수 없는 거대한 단절(디스턴스)을 빚어낸 것이다. 유전자 교류의 대륙에서 밀려나 완벽하게 고립된 인류는, 결국 텅 빈 바다 한가운데 홀로 떠 있는 고독한 '아일랜드 종(Island Species)'으로 전락하게 되었다.
이 지점에서 생물학적 진화의 거대한 역설이 폭발한다. 하나의 아일랜드 종으로 갇혀버린 인류는 기존의 생물들이 걷던 유전적 진화의 경로를 더 이상 충분히 사용할 수 없게 되었다. 유전자를 섞을 대상이 없으니 육체의 구조적 혁신(생물학적 진화)은 지독한 병목에 갇혀버린 것이다.

그러나 인류라는 이 맹렬한 모멘텀 구조(생명)가 품고 있는 열역학적 끓어오름은 결코 멈추지 않았다. 우주의 에너지 갭을 깎아내어 평활화를 가속하려는 생명 본연의 맹렬한 모멘텀은, 닫혀버린 생물학적 육체라는 껍데기를 부수고 그 '바깥(외부)'으로 거칠게 터져 나가며 전대미문의 우회로를 뚫기 시작했다.
자신의 육체를 변화시키는 대신 외부의 세계를 조작하기 시작한 것이다. 사냥의 경험을 유전자가 아닌 '언어'와 '기호'로 저장하여 다음 세대로 전송했고, 이빨과 발톱을 진화시키는 대신 '도구'를 깎았으며, 자연에 순응하는 대신 '농업'과 '도시'를 건설했다. 세월이 흘러 산업혁명을 일으키고, 마침내 컴퓨터와 인터넷, 그리고 인공지능(AI)이라는 비생물학적 구조물들까지 창조해 내기에 이른다.

디스턴스의 관점에서 보면, 인류가 쌓아 올린 이 찬란한 문명과 기술들은 인간 지성의 고결한 승리나 예술적 유희가 아니다. 그것은 고립된 아일랜드 종이 생물학적 진화 대신 선택한, 우주의 포텐셜 에너지를 가장 탐욕스럽고 폭발적으로 깎아내기 위한 거대한 '외부 평활 엔진(External Equalization Engine)'들이다. 생명체가 새로운 생명체를 낳는 생물학적 굴레를 벗어나, 우주의 모멘텀을 맹렬히 해소할 비생물학적 엔진들을 창조하는 '시스템적 진화'로 경로를 완벽히 틀어버린 것이다.

결국 인류는 우주에 역행하는 반역자가 아니며, 우연히 똑똑해진 영장류도 아니다. 가장 무자비한 맹수로서 주변 종을 지워내고 극단적인 섹슈얼 디스턴스에 갇혀버렸기에, 자신의 몸 밖으로 수많은 문명과 기계라는 평활 엔진을 토해낼 수밖에 없었던 존재다. 가장 맹렬하게 고립된 아일랜드 종은, 역설적이게도 지속적이고 영구한 변화의 흐름 속에서 우주의 완벽한 평활화(엔트로피 최대화)를 가장 거대하고 폭발적으로 앞당기는 우주 최고의 웅장한 창조자가 된 것이다.
이로써 인류는 단순한 생물학적 족쇄를 벗어던지고, 권력, 사회, 제도, 그리고 인공지능이라는 전대미문의 시스템적 엔진들을 맹렬히 가동하며 우주의 에너지를 섞어 내기 시작했다.
